\section{Native Klein Generators and Permutations}

This appendix records the exact native permutations underlying the
Klein register. All permutations act on the fixed sixty-vertex
ordering used by the Project 42 certificates.

\subsection{Permutation convention}

A permutation
\[
g\in\operatorname{Sym}(60)
\]
is stored in image-list notation:
\[
[g(0),g(1),\ldots,g(59)].
\]

Permutation multiplication is function composition:
\[
(gh)(v)=g(h(v)).
\]

The identity permutation is
\[
[0,1,\ldots,59].
\]

All subgroup, orbit, conjugation, and stabilizer calculations in this
paper use this convention.

\subsection{The native involution \(a\)}

The involution \(a\) is reconstructed from the thirty canonical
two-vertex fibers of the quotient
\[
G_{60}\longrightarrow G_{30}.
\]

For every canonical fiber
\[
\{u,v\},
\]
the permutation exchanges the two vertices:
\[
a(u)=v,
\qquad
a(v)=u.
\]

The certificate verifies
\[
a^2=1,
\]
\[
a(v)\neq v
\]
for every native vertex \(v\), and
\[
a(E(G_{60}))=E(G_{60}).
\]

The exact image list is stored in
\[
\texttt{native\_g60\_v4\_recovery\_052.json}.
\]

\subsection{The native involution \(b\)}

The second native involution \(b\) is imported from the hash-locked
native generator bundle and verified independently against the exact
edge set.

It satisfies
\[
b^2=1,
\]
has no fixed points, and preserves every native edge.

Its exact image list is stored in
\[
\texttt{native\_g60\_v4\_recovery\_052.json}.
\]

\subsection{The product \(ab\)}

The third nonidentity native Klein element is
\[
ab=a\circ b=b\circ a.
\]

The certificate verifies
\[
(ab)^2=1,
\]
that \(ab\) is fixed-point-free, and that it preserves the exact native
edge set.

Thus
\[
V_4^{\mathrm{nat}}
=
\{1,a,b,ab\}.
\]

\subsection{Group table}

The native Klein multiplication table is

\[
\begin{array}{c|cccc}
 & 1 & a & b & ab\\
\hline
1  & 1  & a  & b  & ab\\
a  & a  & 1  & ab & b\\
b  & b  & ab & 1  & a\\
ab & ab & b  & a  & 1
\end{array}
\]

Every nonidentity element has order two, and the group is abelian.

\subsection{Free action}

For every native vertex \(v\), the orbit
\[
\{v,a(v),b(v),ab(v)\}
\]
contains four distinct vertices.

The fifteen orbits form a partition
\[
V(G_{60})
=
\bigsqcup_{i=0}^{14}C_i.
\]

No nonidentity native Klein element fixes a vertex.

\subsection{Conjugation signatures}

For each
\[
g\in\operatorname{Aut}(G_{60}),
\]
the certificate computes
\[
gag^{-1},
\qquad
gbg^{-1},
\qquad
g(ab)g^{-1}.
\]

Exactly two signatures occur:

\[
(a,b,ab)
\]
for \(240\) automorphisms, and
\[
(a,ab,b)
\]
for \(240\) automorphisms.

The second signature occurs exactly when the induced five-point
permutation is odd.

\subsection{Machine-readable sources}

The authoritative native Klein sources are:

\begin{enumerate}
    \item
    \texttt{artifacts/json/native\_g60\_reconstructed\_deck\_involution\_a\_050.json};

    \item
    \texttt{artifacts/json/native\_g60\_v4\_candidate\_recovery\_051.json};

    \item
    \texttt{artifacts/json/native\_g60\_v4\_recovery\_052.json};

    \item
    \texttt{artifacts/json/native\_g60\_character\_exchange\_c2\_theorem\_056.json};

    \item
    \texttt{artifacts/json/native\_g60\_s5\_extension\_splitting\_audit\_079.json}.
\end{enumerate}

The full image lists are retained in the machine-readable certificates
rather than duplicated in the printed manuscript. This keeps the
appendix readable while preserving exact reproducibility.
